\section{安全与基础设施}

\subsection{日志与可观测性}

在部署 reward learned policy 之前，必须保证完整的日志和可观测性链路。关键维度包括：

\begin{itemize}
    \item \textbf{Reward Trace}：记录 reward model 打分、baseline policy 预测、KL penalty
    \item \textbf{Action Trace}：记录 policy 选择的动作、概率分布、temperature
    \item \textbf{Safety Signals}：检测不当内容、hallucination、重复行为
\end{itemize}

\begin{warningbox}{没有日志，结果就无法纠错}
实际系统中最危险的 reward hacking 不是模型输出歪了，而是因为团队根本看不到哪里“破产”。通用做法是把 reward model score 写进 trace table，并在达不到阈值时自动 roll back。
\end{warningbox}

\subsection{从仿真到现实}

Reward learning 训练常在仿真环境里完成，部署到现实世界前需做“sim-to-real 补充”：

\begin{itemize}
    \item \textbf{Domain randomization}：在训练中加入光照、噪声、物理扰动
    \item \textbf{Reward smoothing}：在 reward signal 上加入dropout或混合真实人类标签，防止 overfit
    \item \textbf{Distillation+Calibration}：把 large policy 蒸馏到小模型，同时校准 reward model 以抵消 distribution shift
\end{itemize}

\begin{knowledgebox}{Sim-to-Real 案例}
在机器人 grasping 任务里，reward learned policy 在仿真超越 95%，但直接部署到真实抓手会下沉至 45%。解决方式是把 sim 收集的 reward signal 与少量真实 human data 混合，并加入形变扰动进行 fine-tuning。
\end{knowledgebox}

\subsection{本章小结}

安全部署要求完整的日志体系、自动审查与 sim-to-real 校准。Reward model 的每一步都应该可追踪、可回滚。
